\documentclass[10pt,a4paper]{amsart}
\usepackage{amsmath,amssymb,amsthm,graphicx}
\usepackage{verbatim}
\newenvironment{reference}{\comment}{\endcomment}
%\newenvironment{reference}{}{}
\newenvironment{slogan}{\comment}{\endcomment}
\newenvironment{history}{\comment}{\endcomment}

% For commutative diagrams we use Xy-pic
\usepackage[all]{xy}

% We use 2cell for 2-commutative diagrams.
\xyoption{2cell}
% Theorem environments.
%
\theoremstyle{plain}
\newtheorem{theorem}[subsection]{Theorem}
\newtheorem{proposition}[subsection]{Proposition}
\newtheorem{lemma}[subsection]{Lemma}

\theoremstyle{definition}
\newtheorem{definition}[subsection]{Definition}
\newtheorem{example}[subsection]{Example}
\newtheorem{exercise}[subsection]{Exercise}
\newtheorem{situation}[subsection]{Situation}

\theoremstyle{remark}
\newtheorem{remark}[subsection]{Remark}
\newtheorem{remarks}[subsection]{Remarks}

\numberwithin{equation}{subsection}

% Macros
%
\def\lim{\mathop{\mathrm{lim}}\nolimits}
\def\colim{\mathop{\mathrm{colim}}\nolimits}
\def\Spec{\mathop{\mathrm{Spec}}}
\def\Hom{\mathop{\mathrm{Hom}}\nolimits}
\def\Ext{\mathop{\mathrm{Ext}}\nolimits}
\def\SheafHom{\mathop{\mathcal{H}\!\mathit{om}}\nolimits}
\def\SheafExt{\mathop{\mathcal{E}\!\mathit{xt}}\nolimits}
\def\Sch{\mathit{Sch}}
\def\Mor{\mathop{\mathrm{Mor}}\nolimits}
\def\Ob{\mathop{\mathrm{Ob}}\nolimits}
\def\Sh{\mathop{\mathit{Sh}}\nolimits}
\def\NL{\mathop{N\!L}\nolimits}
\def\CH{\mathop{\mathrm{CH}}\nolimits}
\def\proetale{{pro\text{-}\acute{e}tale}}
\def\etale{{\acute{e}tale}}
\def\QCoh{\mathit{QCoh}}
\def\Ker{\mathop{\mathrm{Ker}}}
\def\Im{\mathop{\mathrm{Im}}}
\def\Coker{\mathop{\mathrm{Coker}}}
\def\Coim{\mathop{\mathrm{Coim}}}

% Boxtimes
%
\DeclareMathSymbol{\boxtimes}{\mathbin}{AMSa}{"02}

%
% Macros for moduli stacks/spaces
%
\def\QCohstack{\mathcal{QC}\!\mathit{oh}}
\def\Cohstack{\mathcal{C}\!\mathit{oh}}
\def\Spacesstack{\mathcal{S}\!\mathit{paces}}
\def\Quotfunctor{\mathrm{Quot}}
\def\Hilbfunctor{\mathrm{Hilb}}
\def\Curvesstack{\mathcal{C}\!\mathit{urves}}
\def\Polarizedstack{\mathcal{P}\!\mathit{olarized}}
\def\Complexesstack{\mathcal{C}\!\mathit{omplexes}}
% \Pic is the operator that assigns to X its picard group, usage \Pic(X)
% \Picardstack_{X/B} denotes the Picard stack of X over B
% \Picardfunctor_{X/B} denotes the Picard functor of X over B
\def\Pic{\mathop{\mathrm{Pic}}\nolimits}
\def\Picardstack{\mathcal{P}\!\mathit{ic}}
\def\Picardfunctor{\mathrm{Pic}}
\def\Deformationcategory{\mathcal{D}\!\mathit{ef}}
\title{The localization of a ring at the points specializing to a locally closed subscheme of the form D(f) cap V(I)}
\author{The Stacks Project Authors}
\date{stacks-project:proetale.tex:local-snapshot-20261012; GNU FDL 1.2 or later}
\begin{document}
\maketitle
\noindent\textit{Editorial note.} All delivered mathematical text is the author's own, copied from the pinned Stacks Project chapter file proetale.tex; no statement or proof step was rewritten, reordered or removed.\par\medskip
\noindent\textit{Editorial note.} The counted proof is the author's complete argument: the identification of the multiplicative subset through the ring of sections on Z, the two-sided characterization of the primes surviving in the localization, and the formal derivation of the dependence, the closed subscheme statement and the functoriality.\par\medskip
\noindent\textit{Editorial note.} The label lemma-localization names the counted statement, and the proof environment immediately following it in the source is delivered verbatim as its proof.\par\medskip
\noindent\textit{Editorial note.} The counted statement is carried by the author's own lemma environment in the chapter source.\par\medskip
\noindent\textit{Editorial note.} Pre-open gate: python operations/agent-authoring/tools/precard\_gate.py --source proetale.tex --label lemma-localization --start 709 --end 748 --local-id 3596 returned RESULT=OK (L3596 free, occupancy FREE, public interval clear). Candidate taken from sources/stacks-proetale/candidates-halfA.json.\par\medskip
\noindent\textit{Editorial note.} The statement and proof are copied from the pinned chapter file; the cited inputs are kept as external pointers with their canonical Stacks tags.\par\medskip
\noindent\textit{Allowed standard prerequisites (editorial).}\begin{itemize}
\item Description of sections of the structure sheaf on a basic open of a closed subscheme. Algebra and Schemes, the identification of the global sections of the structure sheaf of D(f) cap V(I) with the localization of A/I at f. Supplies the description of the ring whose invertible elements define the multiplicative subset S.
\item Localization and the spectrum of a multiplicative subset. Algebra, the standard description of the primes of a localization as the primes of A disjoint from the multiplicative subset. Supplies the two inclusions characterizing which primes lie in the spectrum of S\textasciicircum{}\{-1\}A.
\end{itemize}\medskip
\begin{lemma}
\label{lemma-localization}
Let $A$ be a ring. Set $X = \Spec(A)$. Let $Z \subset X$ be a locally closed
subscheme which is of the form $D(f) \cap V(I)$ for some $f \in A$ and
ideal $I \subset A$. Then
\begin{enumerate}
\item there exists a multiplicative subset $S \subset A$ such that
$\Spec(S^{-1}A)$ maps by a homeomorphism to the set of points of $X$
specializing to $Z$,
\item the $A$-algebra $A_Z^\sim = S^{-1}A$ depends only on
the underlying locally closed subset $Z \subset X$,
\item $Z$ is a closed subscheme of $\Spec(A_Z^\sim)$,
\end{enumerate}
If $A \to A'$ is a ring map and $Z' \subset X' = \Spec(A')$ is a
locally closed subscheme of the same form which maps into $Z$,
then there is a unique $A$-algebra map
$A_Z^\sim \to (A')_{Z'}^\sim$.
\end{lemma}

\begin{proof}
Let $S \subset A$ be the multiplicative set of elements which map
to invertible elements of $\Gamma(Z, \mathcal{O}_Z) = (A/I)_f$.
If $\mathfrak p$ is a prime of $A$ which does not specialize to $Z$,
then $\mathfrak p$ generates the unit ideal in $(A/I)_f$. Hence
we can write $f^n =  g + h$ for some $n \geq 0$, $g \in \mathfrak p$,
$h \in I$. Then $g \in S$ and we see that $\mathfrak p$ is not in
the spectrum of $S^{-1}A$. Conversely, if $\mathfrak p$ does specialize
to $Z$, say $\mathfrak p \subset \mathfrak q \supset I$ with
$f \not \in \mathfrak q$, then we see that $S^{-1}A$ maps to
$A_\mathfrak q$ and hence $\mathfrak p$ is in the spectrum of $S^{-1}A$.
This proves (1).

\medskip\noindent
The isomorphism class of the localization $S^{-1}A$ depends only
on the corresponding subset $\Spec(S^{-1}A) \subset \Spec(A)$, whence
(2) holds. By construction $S^{-1}A$ maps surjectively onto
$(A/I)_f$, hence (3). The final statement follows as the multiplicative subset
$S' \subset A'$ corresponding to $Z'$ contains the image of the
multiplicative subset $S$.
\end{proof}

\end{document}
